\section{The trace determinants on a fixed contour}\label{sec:local}
This section proves the local partition formula and makes its finite-precision requirements explicit. All notation is as in Theorem~\ref{thm:local-orders}. By compactness of the boundary, $H(t,z)$ remains nonzero there for all sufficiently small complex $t$. The argument principle fixes the enclosed degree $M$, and proves that every $S_j$ is holomorphic in $t$.

\subsection{Hermite splitting and the leading coefficient}
At a zero $c$ of multiplicity $m$, write $H(0,c+y)=a_my^m+O(y^{m+1})$, where $a_m\ne0$. The heat equation recursively determines every time derivative by $\partial_t^nH=(-1)^n\partial_z^{2n}H$. Consequently, on putting $t=u^2/2$ and $z=c+uy$,
\[
 \frac{H(u^2/2,c+uy)}{a_mu^m}
 =\sum_{\ell=0}^{\lfloor m/2\rfloor}
 \frac{(-1)^\ell m!}{2^\ell\ell!(m-2\ell)!}y^{m-2\ell}+O(u)
 =\He_m(y)+O(u).
\]
The series is convergent locally uniformly in $y$. To see the divisibility, the coefficient of $u^{2n+j}$ comes from $\partial_z^{2n+j}H(0,c)$ and vanishes for $2n+j<m$; joint holomorphy supplies convergence. The monic probabilists' Hermite polynomial $\He_m$ has $m$ distinct real roots $h_{m,j}$. Its three-term recurrence is the characteristic-polynomial recurrence of the symmetric tridiagonal matrix with off-diagonals $\sqrt1,\ldots,\sqrt{m-1}$. That matrix has simple eigenvalues: an eigenvector is determined by its first coordinate through the tridiagonal equations, and a zero first coordinate forces all coordinates to vanish. The implicit function theorem now gives all the branches
\[
 z_{ij}(t)=c_i+\sqrt{2t}\,h_{m_i,j}+O(t).
\]
For small positive $t$ they are real, by conjugation and uniqueness.

For $0\le a\le m$ define
\[
 C(m,0)=1,\qquad
 C(m,a)=m^a\prod_{\ell=1}^{a-1}(m-\ell)^{a-\ell}\quad(a\ge1).
\]
\begin{proposition}[The exact leading coefficient]\label{prop:leading-coefficient}
With $v=v_k$ from Theorem~\ref{thm:local-orders}, the coefficient $c_k=[t^v]D_k(t)$ is
\[
 c_k=2^v\!\sum_{\substack{0\le a_i\le m_i,\ \sum_i a_i=k\\
                  \sum_i\binom{a_i}{2}=v}}
       \prod_i C(m_i,a_i)\prod_{i<j}(c_j-c_i)^{2a_i a_j}>0.
\]
\end{proposition}
\begin{proof}
For the $M$ enclosed roots $z_1,\ldots,z_M$, repeated with multiplicity, let $V_{u,j}=z_j^u$ for $0\le u<k$. Then the moment matrix is $VV^{\mathsf T}$. Cauchy--Binet gives
\[
 D_k=\sum_{|I|=k}\prod_{\substack{p<q\\p,q\in I}}(z_q-z_p)^2,
\]
the classical subdiscriminant identity; see Roy \cite[Proposition~1]{roy}. A selection containing $a_i$ roots from the $i$th center has order $e(a)=\sum_i\binom{a_i}{2}$. Each internal squared difference contributes one power of $t$; cross-center differences have nonzero limits. The successive costs at center $i$ are $0,1,\ldots,m_i-1$. Choosing the $k$ least costs respects each center's preceding costs, so the minimum of $e(a)$ is $v_k$. Every term attaining the minimum has a positive leading coefficient. This proves the order formula without any cancellation.

It remains to evaluate the sum of the internal squared Vandermonde factors on $a$ Hermite roots. Let $J_m$ be the tridiagonal matrix just described. Cramer's rule at its last vertex gives the resolvent entry $\He_{m-1}(z)/\He_m(z)$. Differentiating the Hermite recurrence, or its displayed coefficients, gives $\He_m'=m\He_{m-1}$. The spectral measure at that vertex therefore puts mass $1/m$ at each root. Reverse the vertices. The monic orthogonal-polynomial recurrence now has coefficients $m-1,m-2,\ldots,1$. The squared norm of its degree-$j$ monic polynomial is $\prod_{\ell=1}^j(m-\ell)$: apply that polynomial in $J_m$ to the first basis vector, obtaining the product of the first $j$ off-diagonal entries times the $(j+1)$st basis vector. Changing from monomials to these monic polynomials has determinant one. Thus the degree-$a$ moment determinant for the counting measure is
\[
 m^a\prod_{j=0}^{a-1}\prod_{\ell=1}^j(m-\ell)=C(m,a).
\]
Substitution into Cauchy--Binet proves the coefficient formula. Finally, the multiplicity of the integer $h$ among the costs is exactly the number of centers with $m_i\ge h+1$, proving the reconstruction assertion.
\end{proof}

At zero time the moment form is $\sum_i m_iP(c_i)Q(c_i)$, of rank $r$. In particular
\[
 D_r(0)=\prod_i m_i\prod_{i<j}(c_j-c_i)^2>0,
 \qquad D_k(0)=0\quad(k>r).
\]
\begin{proposition}[The first negative local minor]\label{prop:first-negative}
Set $Q(z)=\prod_i(z-c_i)$. For $r<M$,
\[
 D_{r+1}'(0)=2D_r(0)\sum_i m_i(m_i-1)Q'(c_i)^2>0.
\]
Hence $D_{r+1}(-a)<0$ for all sufficiently small $a>0$.
\end{proposition}
\begin{proof}
Replace the degree-$r$ monomial in the moment basis by the monic polynomial $Q$. This basis change has determinant one. Its last row and column vanish at $t=0$, so differentiating the determinant leaves $D_r(0)$ times the derivative of the moment of $Q^2$. For a polynomial $P$, differentiate its contour moment, use $H_t=-H_{zz}$, and integrate by parts on the closed contour to obtain
\[
 \left.\frac d{dt}\sum_zP(z)\right|_{t=0}
 =\sum_i\operatorname*{Res}_{z=c_i}
      P'(z)\frac{H_{zz}(0,z)}{H(0,z)}.
\]
Writing $H(0,z)=(z-c_i)^{m_i}g_i(z)$, the double-pole coefficient is $m_i(m_i-1)$ and the simple-pole coefficient is $2m_i g_i'(c_i)/g_i(c_i)$. For $P=Q^2$, $P'(c_i)=0$ and $P''(c_i)=2Q'(c_i)^2$. The residue is therefore the summand asserted. At least one $m_i$ exceeds one.
\end{proof}

\subsection{Two partitions worked out}
For $P(z)=z^3(z-1)$, its heat evolution is
\[
 e^{-tD^2}P=z^4-z^3-12tz^2+6tz+12t^2.
\]
Newton's identities give
\[
 D_1=4,\quad D_2=3+96t,\quad
 D_3=36t+864t^2+9216t^3,
\]
and the leading term of $D_4$ is $864t^3$. The costs are $0,0,1,2$, giving orders $(0,0,1,3)$. For $P(z)=z^2(z-1)^2$, the costs are $0,0,1,1$, giving orders $(0,0,1,2)$; its leading coefficients for $D_1,D_2,D_3,D_4$ are respectively $4,4,32,64$. The complete determinant polynomials are computed by exact rational arithmetic in the workbench. The implementations compare all coefficients, as well as orders and leading terms.

The closeness of different centers affects the signal. For $P(z)=z^2-\delta^2$,
\[
 D_2(t)=4\delta^2+8t,
\]
so a simple endpoint pair gives a negative minor at times $t<-\delta^2/2$. For $P(z)=z^2(z-\delta)$,
\[
 D_3(t)=8t(\delta^4+18\delta^2t+108t^2).
\]
Its detecting derivative is $8\delta^4$, tending to zero with the distance to the other center. These exact models explain why a negative observation at a prescribed nonzero time needs endpoint control before it can certify multiplicity at time zero.

\subsection{Explicit contour and coefficient errors}
Choose a complex time radius $R>0$ on which the boundary is zero-free. Let $\ell$ be the boundary length and put
\[
 Z=\max(1,\max_{z\in\partial U}|z|),\quad
 A=\frac{\ell}{2\pi}\max_{\substack{|t|\le R\\z\in\partial U}}
       \left|\frac{H_z(t,z)}{H(t,z)}\right|,
 \qquad B_k=k^{k/2}A^kZ^{k(k-1)}.
\]
These are finite numbers for this contour. The boundary integral gives $|S_j(t)|\le AZ^j$. Divide row $u$ and column $v$ of the moment matrix by $Z^u$ and $Z^v$ and apply Hadamard's determinant inequality. Thus
\begin{equation}\label{eq:det-bound}
 |D_k(t)|\le B_k\qquad(|t|\le R).
\end{equation}
For $r>1$ put $\delta=\min_{i<j}|c_i-c_j|>0$; for $r=1$ put $\delta=1$. Proposition~\ref{prop:leading-coefficient} immediately gives
\begin{equation}\label{eq:lead-lower}
 c_k\ge 2^{v_k}\delta^{k(k-1)-2v_k}.
\end{equation}
Indeed, every factor $C(m_i,a_i)$ is a positive integer, at least one minimizing selection exists, and the cross-center exponent is $k(k-1)-2v_k$.

\begin{proposition}[A coefficient certificate]\label{prop:coefficient-certificate}
For $0<\varrho<R$, take $N$ equally spaced points $t_l=\varrho e^{2\pi il/N}$. Let $\widetilde D_l$ approximate $D_k(t_l)$ with absolute error at most $\varepsilon$. For $0\le j<N$ define
\[
 \widetilde d_j=\frac1{N\varrho^j}\sum_{l=0}^{N-1}
           \widetilde D_l e^{-2\pi ilj/N}.
\]
Then
\[
 |\widetilde d_j-[t^j]D_k(t)|
 \le \varepsilon\varrho^{-j}
   +B_kR^{-j}\frac{(\varrho/R)^N}{1-(\varrho/R)^N}.
\]
Furthermore, for $|t|\le R/2$,
\[
 |D_k(t)-c_kt^{v_k}|
 \le 2B_k|t|^{v_k+1}/R^{v_k+1}.
\]
\end{proposition}
\begin{proof}
Write $D_k(t)=\sum_{n\ge0}d_nt^n$. Cauchy's estimate from \eqref{eq:det-bound} is $|d_n|\le B_kR^{-n}$. The discrete Fourier sum of exact samples equals $d_j+\sum_{h\ge1}d_{j+hN}\varrho^{hN}$, since $j<N$. Summing the geometric majorant proves the aliasing bound, and averaging the sample errors proves the first term. The lower Taylor coefficients vanish through degree $v_k-1$ by Theorem~\ref{thm:local-orders}. Summing the coefficient bounds above degree $v_k$ proves the remainder estimate.
\end{proof}

For example, a lower bound $c\le c_k$ and the last estimate certify the sign of $D_k(-a)$ for
\[
 0<a<\min\{R/2,\ cR^{v_k+1}/(2B_k)\}.
\]
Exact recovery of an order from approximate coefficients uses the uniform lower bound on its first nonzero coefficient in \eqref{eq:lead-lower}, scanning $0\le j\le\binom{k}{2}$ with an error below one third of the minimum of those lower bounds. Coefficients preceding the first nonzero coefficient have absolute value at most that error; the first nonzero one has absolute value at least twice that error.

There is also a direct route to the time disk. Let $\eta=\min_{\partial U}|H(0,z)|>0$ and let $V$ bound $|H_t|$ on $|t|\le R$, $z\in\partial U$. The estimate $RV<\eta$ proves boundary nonvanishing by integrating in $t$. A bound for $H_z$ then bounds $A$ above. This identifies the precise analytic estimates required as the contour grows.
